% 本文件由 LaTeX 排版 Agent 根据有效 translation.json 自动生成。
% author=Augustine of Hippo; work_id=1403; title=论两个灵魂

\chapter{论两个灵魂}

% structure: p0001 source=h1 target=omitted reason=page-title-already-rendered-by-parent

\Emph{反摩尼教人（公元391年）}\par

% structure: p0003 source=h2 target=section reason=relative-heading-rank; base=h2

\section{第一章——通过何种推理途径证明摩尼教关于两个灵魂（其中一个非出自天主）的谬误。每个灵魂，既然是一种生命，其存在只能来自生命的根源——天主}

1. 藉着天主仁慈的助佑，在粉碎并摆脱了摩尼教人的罗网之后，我终于被恢复到了天主教会的怀抱，如今我至少愿意反思并哀悼我近来的不幸。因为有许多事我本应去做，以防止那从孩提时代就健康地播种在我心中的最真实宗教的种子，被虚假和诡诈之人的错误和欺骗从我心中拔除。因为，首先，如果我曾以祈祷和虔敬的心，清醒而勤勉地思考过那两类灵魂（摩尼教人赋予它们如此不同的本性和属性，以致他们希望其中之一被视为出于天主的本体，却甚至不愿意承认天主是另一类的作者），也许对于一个专心学习的人来说，我就会明白：没有任何生命，仅凭其是生命并且只要是生命，就不归属于那至高的生命根源和开端，而这根源和开端我们必得承认就是至高的、唯一的、真实的天主。因此，我们没有理由不承认：摩尼教人称之为邪恶的那些灵魂，要么缺乏生命而非灵魂，且不会有任何肯定或否定的行动，不追求也不逃避任何事物；要么，如果它们活着以致能成为灵魂，并如摩尼教人所设想的那样行动，那么它们除非藉着生命就无法活着，而如果这是确定的事实（正如基督所说：\BibleQuote{我是生命}\BibleRef{若 14:6}），那么所有灵魂，既然它们除非藉着活着就不能成为灵魂，就都是由基督、即由生命所创造和形成的。\par

% structure: p0005 source=h2 target=section reason=relative-heading-rank; base=h2

\section{第二章——如果感官所感知的光以天主为作者（如摩尼教人所承认的那样），那么唯独理智所感知的灵魂更是如此}

2. 但是，如果当时我的思想不能承受和坚持有关生命和分享生命的问题（这确实是一个大问题，需要在有学问的人中进行平静的讨论），那么我也许能够发现一个对所有反省自身的人（无需研究个别部分）完全明显的事实：即，我们据说知道和理解的一切事物，我们或通过身体感官，或通过心智活动来理解。五种身体感官通常被列举为视觉、听觉、嗅觉、味觉、触觉，而理智远远比所有这些更高贵、更卓越。谁会如此忘恩负义和不虔敬，而不承认这一点呢？一旦这点得以确立和确认，我们就会看到，由此可以推出：凡通过触觉或视觉或任何身体方式所感知的事物，都比我们理智理解的事物低劣，其低劣程度与感官低于理智的程度相当。因此，既然一切生命，从而每一个灵魂，都不能被任何身体感官所感知，而只能被理智所感知；而同时，摩尼教人自己也说，那边太阳、月亮以及每一颗被这些尘世之眼所见的光体都必须归于真实而良善的天主，那么，我们观察到属于身体的竟归于天主，这是极端的疯狂；而另一方面，我们认为我们不仅用理智，而且用最高形式的思想（即理性和智性）所接受的东西，即生命（无论它被称为什么，仍然是生命），竟然被剥夺和丧失同一作者的天主，这更是极端疯狂。因为，如果我在呼求天主后自问：生命是什么？它对所有身体感官是多么不可思议，多么绝对无形体，难道我不能回答吗？或者，难道摩尼教人不是也承认他们所憎恶的灵魂不仅活着，而且不朽地活着？并且基督所说的\BibleQuote{任凭死人去埋葬他们的死人}\BibleRef{玛 8:22}并不是针对那些根本没有活着的人，而是针对罪人——这是不朽灵魂唯一的死亡；正如\Person{保禄}所写：\BibleQuote{那惯于享乐的寡妇，虽生犹死}\BibleRef{弟前 5:6}，他说她同时是死的和活的。因此，我本应关注的不是有罪的灵魂生活在何等大的污染中，而仅仅是它活着这一事实本身。但是，如果我只能通过智力活动来感知，我相信我会想到：任何心灵无论比我们通过这些眼睛所看见的光优越多少，我们也应该给予理智相对于眼睛本身的优越性。\par

% structure: p0007 source=h2 target=section reason=relative-heading-rank; base=h2

\section{第三章——如何证明每一个身体也来自天主。摩尼教人称之为邪恶的灵魂比光更优越}

他们确实断言光来自基督的圣父；那么，难道我会怀疑每个灵魂都来自祂吗？但是，甚至在那时，像我这样如此没有经验和如此年轻的人，我也会毫不疑惑地认为不仅灵魂，而且身体，甚至一切事物，都来自祂，如果我曾虔敬而谨慎地反思过：形式是什么，或已经形成的是什么，形状是什么，以及被赋予形状的是什么。\par

3. 但現在暫且不談身體，我要哀歎靈魂、自發而生動的運動、行動、生命、不朽；總之，我哀歎我這可憐人竟然相信任何事物能離了天主的良善而擁有這一切屬性；而這些屬性是如此偉大，我卻可悲地忽略不去思考；我認為這對我應是哀歎和哭泣的事。我本該內心深思這些事，與自己討論，將它們轉述給他人，提出問題：認知的力是什麼，既然人身上沒有任何事物能與此卓越相比？而作為人，只要他們是人，就會同意我這一點；我本該詢問：用這些眼睛看見是否就是認知？假如他們回答不是，我本該首先斷定：心智的智力遠低於視覺感覺；然後我本該補充：我們藉由更好的事物所感知的，本身必然被判斷為更好。誰不會同意這一點呢？我本該繼續詢問：他們稱為\Quote{惡}的那個靈魂，是視覺感覺的對象還是心智智力的對象？他們會承認是後者。這一切在我們之間達成共識並確認後，我本該證明由此得出：他們所咒罵的那個靈魂，實在比他們所敬拜的那個光更好，因為前者是心理認知的對象，後者是肉體感覺的對象。但或許他們會在這裡停住，不願跟隨理性的引導；根深蒂固的意見以及長期維護和相信的謬誤有多大的力量啊！但我本該在他們停頓時更加督促他們，不是嚴厲地、不是幼稚地、不是固執地；我本該重複已被承認的事，並表明它們是如何必須被承認的。我本該勸告他們共同商議，以便清楚看出必須向我們否認的是什麼：他們是否認為理智的感知比這些肉體的眼睛更優越是假的，或者藉由心智的卓越所認識到的比藉由卑賤的肉體感覺所認識到的更卓越？他們是否不願意承認，那些他們認為是異類的靈魂，只能藉由理智的感知——即藉由心智的卓越本身——被認識？他們是否想否認太陽和月亮只能藉由這些眼睛才為我們所知？但如果他們回答，這些事物中沒有一樣能被否認，而不至於最荒謬和最無恥，我本該敦促他們不應懷疑：他們宣告值得敬拜的那個光，比他們勸告人們要逃避的那個靈魂更低劣。\par

% structure: p0010 source=h2 target=section reason=relative-heading-rank; base=h2

\section{第四章——連蒼蠅的靈魂也比光更卓越}

4. 在這裡，若他們在困惑中偶然問我，我是否認為連蒼蠅的靈魂都勝過那個光，我本該回答：是的；蒼蠅雖小，不應使我困擾，反而應堅定我認為它是活的。因為要問的是：是什麼使那些微小的肢體生長？是什麼引導如此細小的身體按自然慾望四處移動？是什麼在奔跑時以數字的次序移動它的腳？是什麼在飛翔時調節並振動它的翅膀？無論這是什麼，在如此微小的受造物中，對一個仔細思考的人而言，它如此突出地高聳，以至於勝過任何閃爍在眼前的閃電。\par

% structure: p0012 source=h2 target=section reason=relative-heading-rank; base=h2

\section{第五章——惡習的靈魂，無論多麼該受譴責，如何超越在其種類中可受讚美的光}

當然，沒有人懷疑：任何理智感知的對象，憑藉神聖的法律，在卓越上勝過一切可感覺的對象，因此也勝過這個光。因為，我問，我們藉由思想感知到什麼，如果不是：用心智認識是一回事，經歷身體感覺是另一回事，前者比後者無比地崇高，因此可知的事物必須優先於可感覺的事物，因為理智本身是如此高舉於感官之上？\par

5. 因此我或許也本該知道，以下這一點是明顯的：既然不義、不節制以及其他心智的惡習不是感官的對象，而是理智的對象，那麼這些我們所憎惡並認為可譴責的事物，是如何因為它們是理智的對象，而能超越這個光，無論它在自己的種類中多麼值得讚美。因為，一個良好順服於天主的頭腦深刻意識到：首先，並非每一件我們讚美的事物都優先於我們指責的每一件事物。因為當我們讚美最純淨的鉛時，我並不一定因此將它看得比我所指責的黃金更高。因為每樣事物都必須在其種類中考量。我反對一個不熟悉許多法條的律師，但我如此優先於那個最受讚賞的裁縫，以至於我認為他無比優越。但我讚美裁縫，因為他精通自己的手藝，而我合理地責備律師，因為他不完美地履行其職業的職能。因此我本該發現：在其種類中完美的光，正當地受到讚美；然而因為它被歸入可感覺事物的行列，而這類事物必須讓位於可知事物的行列，所以它必被列於不義和不節制的靈魂之下，因為後者是可知的；雖然我們可以無不公地判斷這些靈魂最該受譴責。因為對於這些靈魂，我們求他們與天主和好，而不是求他們被置於那閃電之上。因此，若有人曾主張這個發光體來自天主，我不會反對；反而我會說：靈魂，即使是惡劣的靈魂，不是就其惡劣而言，而是就其為靈魂而言，必須被承認為天主的受造物。\par

% structure: p0015 source=h2 target=section reason=relative-heading-rank; base=h2

\section{第六章——论恶习本身作为理智把握的对象是否应优于感官对象的光，以及是否应归因于天主为其作者。心灵之恶习及其某些缺陷不应被正当地算作可理解的事物。即便将缺陷本身算作可理解的事物，也绝不应置于可感觉的事物之前。如果光可由天主看见，那么灵魂更是如此，即使有恶习，就它活着而言，它是一个可理解的事物。摩尼教徒所引的相反经文。}

此时，假设他们中有一位谨慎而警觉者，如今更勤学而非固执，曾告诫我说，所探讨的对象不是有恶习的灵魂，而是恶习本身；恶习既非由身体感官所知，却又为人所知，故只能被看作理智把握的对象；倘若这些理智对象超越一切感官对象，那么我们为何不能一致将光归因于天主为其作者，而只有亵渎者才会说天主是恶习的作者？——我本应对这人作出回应，若当时灵机一动，如同善神崇拜者的惯常情形，从高处如闪电般闪现对这一问题的解答，或者解答已预先准备好。若我未曾堪当或无法利用这两种方法中的任何一种，我便应推迟此事，承认所提出的问题难以分辨且艰巨。我应退而自省，俯伏于天主前，大声呻吟，求祂不让我在半途停滞——我本应以确凿的论证前进——求祂不致因一个可疑的问题而被迫要么将可理解之物置于可感觉之物之下并屈服，要么称祂自己为恶习的作者；因为这两种选择都完全充满了虚假和不敬。我绝无可能设想祂会在我这种心态下抛弃我。相反，祂会以祂不可言喻的方式，再三劝诫我去思考那令我如此困扰的心灵恶习是否应算作可理解之物。然而，为能查明——由于我内心之眼的软弱，这因我的罪而合理临到我——我应在可见事物本身中设计某种观照精神事物的舞台；对于这些可见事物，我们绝非有更确定的知识，却有更自信的熟悉感。因此，我应立即追问：什么恰当地属于眼睛的感觉？我应发现那是颜色，而光主宰着颜色。因为这些是其他感官所不触及的；至于物体的运动、大小、距离和形状，虽然它们也能被眼睛感知，但感知这些并非眼睛的独特功能，也属于触觉。由此我应得出结论：那边之光超越其他有形可感觉之物多少，视觉就比其他感官高贵多少。因此，从所有被身体感官感知的事物中选出光，我应努力以此光为阶梯，并必然将我的追问置于其上。我应继续考虑以这种方式可能达成什么，并这样与自己推理：若那边以其光辉夺目、以其光明满足白日的太阳，逐渐在我们眼前黯淡成月亮的模样，我们除了仍存留的光（无论多么暗淡）之外，能用眼睛感知到什么吗？——由于看不到已逝去的（光），我们寻求（现存的）光，并用它来看眼前之物。因此，我们应看不到那个衰退，而只能看到衰退后幸存的光。但既然我们看不到衰退，我们就无法感知它；因为凡由视觉感知的，必定被看见；所以，若那衰退既未被视觉也未被任何其他感官感知，它就不能算作感官对象。因为凡不能被感官感知的，都不是感官对象。现在让我们将这一思考应用于德行，我们最恰当地说心灵因其理智之光而闪耀。再者，从这德行之光的一种衰退——不毁灭灵魂，但使之暗淡——被称为恶习。因此，恶习也绝不能被算作理智感知的对象，正如那光的衰退被正当地排除出感官感知对象的数目一样。然而，灵魂中剩下的东西，即那活着且是灵魂的东西，与那在可见光源中经过任何可想象的衰退后仍应闪耀的东西一样，是理智感知的对象。因此，灵魂，就其是灵魂且分享生命（没有生命，它绝不可能成为灵魂）而言，最正确地应被优先于一切感官感知的对象。所以，说任何灵魂非来自天主——你们夸耀太阳和月亮由祂而来——是极其错误的。\par

7. 但如果现在有人以为，不仅我们通过感官感知的所有事物可称为感官知觉的对象，而且那些我们虽不通过感官感知，却借助身体判断的事物，如通过眼睛判断黑暗，通过耳朵判断寂静——因为我们并非因看见黑暗或听见寂静而知道它们的存在——同样，对于理智知觉的对象，不仅包括那些我们凭借理智光照所看见的事物，如智慧本身，也包括那些我们因这光照本身而避免的事物，如愚昧，我或许可以恰当地称其为理智的黑暗；那么，我就不会就一个词发生争议，而是通过一个简单的划分解决整个问题，并立即向那些留心的人证明，按照天主的真理法则，可理解的\OriginalTerm{自立体}{subsistence}应优先于可感觉的自立体，而不是这些自立体（即使我们愿意称这些为可理解的，那些为可感觉的）的缺失。因此，那些承认这些可见的光体和那些可理解的灵魂都是自立体的人，必然被迫承认并赋予灵魂更崇高的部分；但任何一类\OriginalTerm{缺失}{privation}都不能彼此优先，它们只是缺失，表明非存在，因此相当于\OriginalTerm{否定}{negation}本身。因为当我们说\Quote{这不是金子}和\Quote{这不是美德}时，尽管金子和美德之间存在极大差异，但我们对它们所作的否定之间却没有差异。然而，缺乏美德确实比缺乏金子更糟糕，这是任何心智健全的人都不会怀疑的。谁不知道区别不在于否定本身，而在于它们所附加的事物？因为美德比金子越优越，缺乏美德就比缺乏金子更悲惨。因此，既然可理解的事物优于可感觉的事物，我们理应对可理解事物的缺陷比可感觉事物的缺陷更加反感，不是评价缺陷本身，而是评价有缺陷的事物多多少少的珍贵。由此看来，可理解的光的缺陷远比可感觉的光的缺陷更为悲惨，因为，显然，被认知的生命远比那被看见的光更为珍贵。\par

8. 既然如此，谁敢将太阳、月亮以及星辰中一切光辉的，甚至我们这火和这可见的尘世生命都归于天主，却拒绝承认任何灵魂（灵魂之所以是灵魂，只因它们完美地活着，因为单在这事上，\OriginalTerm{生命}{life}优先于\OriginalTerm{光}{light}）来自天主？而且，说真理的人说：\Quote{事物照耀到什么程度，它就来自天主到什么程度}，那么，大能的天主，如果我说：\Quote{事物活到什么程度，它就来自天主到什么程度}，我岂不是说谎？我恳求你们，不要让\OriginalTerm{理智}{intellect}的盲目和心智的颠倒扩展到如此地步，以致人不能认识这些事。但无论他们的错误和顽固有多大，依赖这些论证并以此武装，我相信，当我将如此思考和讨论过的事向他们提出，并平静地与他们商议时，我会担心，若有人试图将理智或那些（不以缺陷方式）被理智所知觉的事物，从属于甚或比作身体感官，或那些属于身体感官作为知识对象的事物，那么他们中任何人似乎对我有任何影响。这一点确定后，他或任何其他人怎敢否认，那些他认为邪恶的灵魂，既然它们是灵魂，就应算在可理解之物的数目中，并且也不是以\OriginalTerm{缺陷}{defect}方式作为理智知觉的对象？这是基于灵魂只因其活着才成为灵魂的假设。因为如果它们因其缺陷而被理智地知觉为邪恶，因缺乏美德而邪恶，然而它们作为灵魂并不是因缺陷而被知觉，因为它们是因活着而成为灵魂。也不能主张生命的存在是缺陷的原因，因为事物越有缺陷，就越远离生命。\par

9. 因此，既然已经以各种方式显而易见，灵魂不能与那位作者分离，因为那光不与祂分离，那么无论他们现在可能提出什么，我都不会接受，反而要劝告他们，他们应当与我一同选择追随那些主张一切存在者，既因它存在，且无论它以何种程度存在，其存在都来自唯一的天主的人。\par

% structure: p0020 source=h2 target=section reason=relative-heading-rank; base=h2

\section{第七章 恶人如何属于天主又不属于天主}

他们或许会引用福音中那些话反对我：\BibleQuote{所以你们不听，因为你们不是出于天主；} \BibleQuote{你们是出于你们的父亲魔鬼。} 我也应当引用：\BibleQuote{万物是藉着祂造成的，凡受造的，没有一样不是藉着祂造成的}，\BibleRef{若 1:3}以及宗徒的这话：\BibleQuote{只有一个天主，万物都出于祂；并有一位主耶稣基督，万物都是藉着祂而有的}，\BibleRef{格前 8:6}又同一位宗徒说：\BibleQuote{万物都出于祂，藉着祂，在祂内，愿光荣归于祂}。\BibleRef{罗 11:36} 我本应劝诫那些人（如果我确实发现他们是人），不要自以为好像我们已经发现了什么，而应向那些能证明看似不协调的经文之间的一致与和谐的大师请教。 因为当我们读到同一圣经权威说：\BibleQuote{一切都是出于天主}，\BibleRef{格前 11:12}而另一处又说：\BibleQuote{你们不是出于天主}，既然轻率地谴责圣经书卷是错误，谁不会看出需要一位有技巧的教师，他知道这个难题的解决方法呢？从他那必定会，如果他有良好的智力，并且是\BibleQuote{属神的人}，正如因天主感动所说的\BibleRef{格前 2:15}（因为他必然会赞同关于可理解与可感觉的本性的真实论证，这些论证我尽力进行了，也处理了，不，他还会更好、更有说服力地揭示它们）；我们关于这个难题不会再听到别的，只会听到，可能，没有一类灵魂不是从天主获得存在，而对罪人和不信者说\BibleQuote{你们不是出于天主}也是正确的。 因为我们或许，在恳求天主助佑后，本可以容易地看出，活着是一回事，犯罪是另一回事，而且（虽然活在罪中可与义人的生活相比称为死\BibleRef{弟前 5:6}，并且在一个人身上，他同时活着又是罪人），就他活着而言，他是出于天主；就他是罪人而言，他不是出于天主。 在这区分中，我们使用符合我们情感的那个选择；这样，当我们希望强调天主作为创造者的全能时，我们甚至可以对罪人说他们是出于天主。因为我们是在对属于某一类的人说话，是对有生命的人说话，是对有理性的受造物说话，最后——尤其适用于当前问题——是对有生命的存在说话，所有这些都是神圣的职能。 但是当我们的目的是要定罪恶人时，我们正确地对他们说：\BibleQuote{你们不是出于天主}。因为我们是对他们说话，他们背离真理，不信、犯罪、无耻，总之一句话——罪人，这一切无疑都不是出于天主。 因此，如果\Person{基督}对罪人说话，正因为他们犯罪不信祂，而说：\BibleQuote{你们不是出于天主}；而另一方面，又毫无矛盾地说：\BibleQuote{万物是藉着祂造成的}和\BibleQuote{一切都是出于天主}，这有什么奇怪呢？ 因为不相信\Person{基督}、拒绝\Person{基督}的来临、不接受\Person{基督}，确凿无疑是那些不出于天主的灵魂的标志；所以才有这话：\BibleQuote{所以你们不听，因为你们不是出于天主}；那么宗徒那句话怎么是真的呢，那话出现在\BibleBook{若望}{福音}令人难忘的开端：\BibleQuote{祂来到自己的地方，自己的人却没有接受祂}？\BibleRef{若 1:11} 如果他们没有接受祂，他们怎么是自己的？或者因此，因为他们没有接受祂，他们就不是自己的？除非是因为，罪人作为人属于天主，但作为罪人属于魔鬼？ 说\BibleQuote{自己的人没有接受祂}是指本性；而说\BibleQuote{你们不是出于天主}是指意志；因为福音作者在赞扬天主的工程，而\Person{基督}在谴责人的罪恶。\par

% structure: p0022 source=h2 target=section reason=relative-heading-rank; base=h2

\section{第8章。——摩尼教徒探问恶从何而来，并以为藉此问题他们已获胜。让他们首先知道，这是最容易做的，即没有天主，无物能生存。除非藉着对至善——即天主——的认识，至恶不能被认识。}

在此，或许有人会说：罪恶本身从何而来？恶一般从何而来？若来自人，人又从何而来？若来自天使，天使又从何而来？虽然，无论怎样真实且正确地，说这些出于天主，然而对那些不熟练且少有洞察隐秘之事能力的人，恶与罪似乎因此如同链条一般与天主相连。他们以为凭此问题就胜利了，仿佛问就是知——但愿如此，因为那样就没有人比我更知道了。然而，在争辩中，提出重大问题的人，常假扮成伟大的教师，而在所论之事上，他本人比他想要恐吓的对手更为无知。\par

这些人因此以为他们超越常人，因为前者提出了后者无法回答的问题。因此，若在我极为不幸地与他们交往时——并非像我现在这样已有一段时间所处的地位——当我提出这一论证时，他们提出了这些异议，我本应说：我请求你们暂且同意我，这是最容易的事：若没有天主，任何事物都不能发光，那么没有天主，任何事物更不能活着。我们不要固执于如此怪异的见解，竟主张有些灵魂可以离开天主而拥有生命。因为或许可能，你与我一样对这件事——即恶从何而来——一无所知，那么让我们同时或按任何顺序学习，我不在乎什么顺序。因为若没有对完美之善的认识，人就不可能认识完美之恶，岂不是如此？因为我们若常处于黑暗中，就不会认识黑暗。但光的概念不允许它的对立面被隐藏。然而，至善就是那没有比它更高者。但天主是善，没有比祂更高者。因此，天主的至善。因此，让我们一同这样认识天主，这样我们过于匆忙寻求的就不会向我们隐藏。那么，你认为认识天主是一件小事或小功劳吗？因为为我们还有什么其他赏报，除了永生，即认识天主？因为主天主说：\BibleQuote{永生就是：认识你，唯一的真天主，和你所派遣来的耶稣基督。}\BibleRef{若 17:3}因为灵魂虽是不死的，但由于远离对天主的认识被正当地称为它的死亡，当它归向天主时，所获得的永生赏报就是那认识；因此，正如所说，这就是永生。但没有人能归向天主，除非他使自己远离这个世界。对我自己而言，我觉得这既艰难又极其困难，至于对你是否容易，天主自己必会看到。我本倾向于认为对你容易，若非被如下事实所打动：既然我们被命令远离的这个世界是可见的，而宗徒说：\BibleQuote{那看得见的是暂时的，那看不见的才是永远的。}\BibleRef{格后 4:18}你们竟将更多的分量归于这些眼睛的判断而非心灵的判断，竟主张并相信没有不发光的羽毛不是从天主发光，却有从天主不活的灵魂活着。这些话或类似的话，我本应对他们说出或自己思考，因为即使那时，以我全部的心肠恳求天主，并尽可能仔细地查考圣经，我或许也能说出或想到这样的话，只要对我的得救是必要的。\par

% structure: p0025 source=h2 target=section reason=relative-heading-rank; base=h2

\section{第九章——奥古斯丁因与摩尼教徒的亲密往来及其所报告的对无知基督徒的接连胜利而受骗。摩尼教徒也容易从对罪和意志的知识上被驳倒}

但有两件事，尤其容易抓住那不经世事的年龄，通过奇妙的迂回催促着我。一件是亲密往来，突然，借一种虚假的善的外表，如一条蜿蜒的链条多次缠绕我的颈项。另一件是，我在与那些无知的、却仍热切地各尽其力捍卫自己信仰的基督徒辩论时，几乎总是有害地获胜。由于这种频繁的胜利，青年的热情被点燃，并因自身的冲动鲁莽地趋向倔强的大恶。因为这种争辩，在我成为他们中的听者之后，无论我凭自己的天赋（不论如何）或通过阅读他人的著作所能做的，我都极乐意唯独奉献给他们。因此，从他们的言论中，辩论的热情每日增长，从辩论的胜利中，对摩尼教徒的爱恋也增长。结果，无论他们说什么，我仿佛被某种奇怪的疾病所影响，都当成真理而赞同，并非因为我明知其为真理，而是因为我希望它是真理。于是，无论多么缓慢和谨慎，我长久地追随了那些宁愿要光滑的稻草而不要活灵魂的人。\par

12. 就这样，我当时不能分辨感官之物与理智之物，即肉性之物与灵性之物。这不属于年龄、训练、习性，甚至也不属于任何功德；因为这是一件不小的喜乐和福气：若我最终未能把握那在众人判断中本性本身由至高天主的法律所建立的事，那又如何？\par

% structure: p0028 source=h2 target=section reason=relative-heading-rank; base=h2

\section{第十章——罪仅来自意志。各人最了解自己的生命和意志。什么是意志}

因为任凭什么样的人，只要不是疯狂使他们脱离人类的常识，无论他们带着多么大的判断热忱，多么无知，甚至多么迟钝的心智，我想要知道，如果他们被我问道，是否一个人因在他睡着时另一个人用他的手写下了可耻的东西就似乎犯了罪，他们会如何回答？谁会怀疑他们会否认这是罪，并且如此激烈地反对，以至于他们可能因我认为他们适合被问这样的问题而愤怒？在他们被我安抚并恢复平静后，尽我所能，我应请求他们，如果我问他们另一件同样明显、同样众所周知的事情，他们不要见怪。然后我会问，如果某个更强壮的人用另一个并未睡着但清醒的人的手做了一件恶事，而后者其余的肢体被捆绑和束缚，那么因为他知道这事，尽管完全不愿意，他是否应被视为有任何罪过？在这里，所有人都奇怪我竟会问这样的问题，会毫不犹豫地回答，他绝对没有犯任何罪。为什么这样？因为无论谁通过一个无知觉或无法抵抗的人做了任何恶事，后者绝不能被公正地定罪。而正是为什么如此，如果我问这些人的天性，我应当很容易得到想要的答案，通过这样问：假设那个睡着的人已经知道另一个人将用他的手做什么，并预谋地喝了足以阻止他被唤醒的酒，然后去睡觉，以便用誓言欺骗某人。多少睡眠足以证明他的清白呢？除了一个有罪的人，还能称他为什么呢？但如果他也自愿被捆绑，以便借这个借口欺骗某人，那么那些锁链在解除他的罪过方面又有何益处呢？尽管被这些锁链捆绑，他实际上无法抵抗，就像在另一个例子中，睡着的人完全不知道他当时在做什么。因此，难道有任何怀疑两者都应被判为有罪的可能吗？这些被承认后，我就应当论证，罪确实只存在于意志中，因为这一考虑也会帮助我：正义只凭恶的意志就判定那些犯罪的人有罪，尽管他们可能无法实现他们所意愿的。\par

13. 因为谁能说，在提出这些考虑时，我是在谈论晦涩难解的事，而在这些事上，因为能理解的人很少，通常会产生欺诈或炫耀的嫌疑？让可理解的事物与可感觉的事物之间的区分暂时退让一下：不要让我因为用微妙的争论刺激迟钝的心智而受到责难。请允许我知道我活着，请允许我知道我意愿活着。如果人类在这点上一致，如同我们的生命为我们所知，我们的意志也是如此。当我们拥有这种知识时，也无需担心有人会说服我们可能受骗；因为没有人能在自己是否活着或是否一无所愿上受骗。我不认为我提出了任何晦涩的东西，我反而担心有些人会责怪我停留在过于明显的事情上。但让我们思考这些事情的意义。\par

14. 因此，犯罪只通过意志的行使发生。但我们的意志是我们所非常熟悉的；因为如果我不知道意志本身是什么，我也不会知道我在意愿。因此，它被如此定义：意志是心灵的一种运动，无人强迫，为不失去或获得某物。那么为什么我当时不能这样定义它呢？难道很难看出一个不情愿的人与一个情愿的人是相反的，就像左手与右手相反，而不像黑与白相反？因为同一事物不能同时是黑和白。但是，处于两个人之间的人，相对于一个人是在左边，相对于另一个人是在右边。同一个人同时既在右边又在左边，但绝不可能相对于同一个人。确实，同一个心灵可以同时既不情愿又情愿，但它不能相对于同一件事物同时既不情愿又情愿。因为当任何人非情愿地做某事时；如果你问他是否希望做这件事，他说不。同样，如果你问他是否希望不做这件事，他回答说是。这样你会发现他不情愿相对于做，情愿相对于不做，也就是说，同一个心灵同时具有两种态度，但每个态度指向不同的事物。我为什么说这些？因为如果我们再问，他为什么非情愿地做这件事，他会说他是被迫的。因为每个非情愿地做事的人都是被迫的，而每个被迫的人，如果他做了某事，也只是非情愿地做。由此可知，情愿的人不受强迫，即使有人认为他自己是被强迫的。这样，每个情愿地做事的人都不受强迫，而每个不受强迫的人，要么情愿地做，要么根本不做。既然大自然本身在所有人身上宣告了这些事情，从孩童到老人，从文学游戏到智者的宝座，我们都可以毫无荒谬地询问他们，那么为什么我当初没有看到在意志的定义中应该加上\Quote{无人强迫}呢？现在我好像凭借更大的经验极其谨慎地这样做了。但如果这到处都显而易见，并且不是通过教导而是通过自然立刻出现在所有人面前，那么还有什么显得晦涩呢？除非也许有人不知道，当我们渴望某物时，我们意愿，我们的心灵朝它运动，我们要么拥有它要么没有，如果拥有我们就意愿保留它，如果没有，我们就意愿获得它。因此，每个意愿的人，要么意愿不失去某物，要么意愿获得某物。因此，如果所有这些事情都比白天更清楚，正如它们确实是的，而且它们不仅赐予我的理解，而是由真理本身的慷慨赐予全人类，那么为什么我当时甚至连以下定义都不能说呢：意志是心灵的一种运动，无人强迫，为不失去或获得某物？\par

% structure: p0032 source=h2 target=section reason=relative-heading-rank; base=h2

\section{第11章 罪是什么}

有人会说：这会给你提供什么帮助来对抗摩尼派呢？请稍等；让我先也定义罪。每一个心灵都在自身中神性地读到：罪不能脱离意志而存在。因此，罪是保留和追求正义所禁止的、且可以自由避免的事物的意志。尽管如果不自由，它就不是意志。但我更愿意粗略地而非精确地定义。我难道不应该也仔细考察那些晦涩的书籍，从中可以学到，任何人如果要么意愿正义不禁止他意愿的事物，要么不做他不能做的事，就不应受责备或惩罚？难道牧羊人在山上、诗人在剧院中、无学者在社交中、学者在图书馆中、教师在学校中、司铎在圣所中、以及全世界的人类，不都在高唱这些事情吗？但如果任何人，要么不违反正义的禁令行事，要么不做他不能做的事，都不应受责备和谴责，而每一个罪都是应受责备和谴责的，那么当意愿是不义的，且不意愿是自由的时，谁还会怀疑那就是罪呢？因此，这个定义既真实又易于理解，不仅现在，而且当时也能被我这样说：罪是保留或获得正义所禁止的、且可以自由避免的事物的意志？\par

% structure: p0034 source=h2 target=section reason=relative-heading-rank; base=h2

\section{第12章 从所给出的罪和意志的定义，他推翻了摩尼派的整个异端。同样，从对邪恶灵魂的公正谴责，可以推出它们不是因本性而是因意志而邪恶。灵魂因本性是善的，罪过得赦免就是赐予它们。}

16. 现在，让我们看看这些事在哪些方面本会帮助我们。在各种方式上，大有裨益，以致我本来不再需要任何事物；因为它们结束了整个案件。因为一个人若在内心里（在那里它们更清晰明确）省察自己良心的秘密，以及绝对加诸于自然的天主律法，就会承认意志与罪的这两个定义是真实的，并以最少、最简短却最无可反驳的理由毫不犹豫地谴责整个摩尼异端者的异端。这可以这样来思考。他们说有两种灵魂，一种善的，它如此来自天主，以致据说不是祂从任何物质或虚无中造出，而是作为从天主本身实体的一部分而出；另一种恶的，他们相信并试图使别人相信它与天主毫无关系；因此他们主张前者是善的极致，后者是恶的极致，并且这两类曾经是分离的，现在却混合在一起。这个混合的性质与原因我尚未听到；但尽管如此，我本可以问，那恶的一类灵魂在与善的混合之前，是否有任何意志。如果没有，它就没有罪且是无辜的，因此绝不是恶的。但如果它是恶的，以至于虽无意志（如火），但若它接触善，就会侵犯并败坏它，那么相信恶的本性足以改变天主的任何部分，并且至高的善是可朽坏、可侵犯的，这是多么不敬！但如果意志存在，那么毫无疑问，存在一种不受强迫的心灵运动，要么是为了不失去某物，要么是为了获得某物。但这某物要么是善的，要么被认为为善，否则它就不可能被热切地渴望。但在他们主张的混合之前的至高之恶中，从来没有任何善。那么，其中怎能有对善的知识或想法呢？他们是否不渴望自身中的任何事物，却热切地渴望外在的真实之善？这种意志，若热切地渴望至高且真实的善，确实必须被宣称为配得卓越而伟大的赞美。那么，在至高之恶中，这种如此配得如此伟大赞美的心灵运动从何而来？他们是否为了伤害它而寻求它？首先，论证归于同一。因为想要伤害别人的人，是为了自己的某种善而剥夺他人的某种善。因此，在他们之中要么有对善的知识，要么有对善的意见，而这绝不应属于至高之恶。其次，他们如何知道那置于他们之外、他们意图去伤害的善存在呢？如果他们以理智察觉了它，还有什么比这样的心智更卓越呢？善人们全力以赴的不就是为认识那至高而纯正的善吗？因此，如今几乎不赐予少数善人和义人的东西，那恶者（没有善辅助）当时竟能完成吗？但如果那些灵魂带着身体并以肉眼看见了至高之善，那么什么样的舌头、什么样的心灵、什么样的理智足以赞美和宣扬那些眼睛呢？义人们的心灵几乎不能与它们相比！我们在至高之恶中发现了多大的善事！因为如果看见天主是恶的，那么天主就不是善；但天主是善；因此看见天主是善；我不知道还有什么能与此善相比。既然看见任何事物是善的，那么怎能证明能够看见是恶的呢？因此，无论在那些眼睛或那些心灵中有什么使得天主的本质能被它们看见，它都带来了伟大而善的事物，最配得不可言喻的赞美。但如果不是被带来的，而是本身如此且是永恒的，则很难找到比这恶更好的东西了。\par

17. 最后，为使这些灵魂可能没有这些值得赞美的事物（根据摩尼异端者的推理它们被迫拥有），我本应问：天主是否谴责任何灵魂？如果不谴责，就没有赏罚的审判，没有天主眷顾，世界是由偶然而非理智管理，或者更确切地说完全不被管理。因为偶然不能被称为管理。但如果任何敬畏宗教的人都认为这是不敬的，那么要么有些灵魂被谴责，要么没有罪。但如果没有罪，也就没有恶。如果摩尼异端者这样说，他们就会一击杀死自己的异端。因此，他们和我都同意，有些灵魂被天主的法律和审判所谴责。但如果这些灵魂是善的，那是什么正义？如果是恶的，它们是因本性还是因意志而恶？但灵魂决不可能因本性而恶。我们从何处教导这一点？从上述意志和罪的定义。因为说到灵魂，并且说它们是恶的，而又说它们不犯罪，这是疯狂；但说它们没有意志而犯罪，是极大的愚蠢，而认为一个人因没有做他不能做的事而犯罪，属于极端的邪恶和疯狂。因此，无论这些灵魂做什么，如果是因本性而非因意志而行，即为，它们缺乏一种既可做也可不做的自由心灵运动，最后，如果没有任何能力能戒除它们的工作，我们就不能认为罪属于它们。但所有人都承认，恶的灵魂被公正地谴责，而没有犯罪的灵魂被不公正地谴责；因此他们承认那些犯罪的灵魂是恶的。但正如理性所教导的，它们不犯罪。因此，摩尼异端者那种异类的恶灵魂，无论它是什么，都是不存在之物。\par

18. 现在让我们来看看那类善类灵魂，他们又将其高举到如此程度，以至于说它是\Person{天主}的实体。但每个人能认清自己的等级和功德，而不被亵渎的骄傲冲昏头脑，以至相信每当自己经历变化时，那就是至善的实体——而虔诚的理性坚持并教导至善是不可变的——这该多好啊！因为，请看！既然明显灵魂不会因不能成为它们所不能成为的而犯罪，那么这些假设的灵魂，无论它们是什么，根本不犯罪，而且它们根本就不存在；剩下的就是，既然有罪存在，他们除了善类灵魂和\Person{天主}的实体之外找不到任何人可以归咎。但尤其是基督徒的权威使他们感到压力；因为他们从未否认当任何人归向\Person{天主}时，罪得赦免；他们从未说过（正如他们对许多其他段落所说）某个败坏者将此插入圣经。那么，罪归给谁呢？如果归给异类恶灵魂，这些灵魂也可以变好，可以与\Person{基督}一同承受\Person{天主}的国。他们（摩尼教徒）否认这一点，那么除了那些他们主张是\Person{天主}实体的灵魂之外，就没有其他类灵魂了。剩下的就是，他们承认不仅后者犯罪，而且只有后者犯罪。但我不争论他们是否单独犯罪；然而他们确实犯罪。但他们是因与恶混杂而被强迫犯罪吗？如果是被强迫以至于没有抵抗的能力，那么他们就没有犯罪。如果抵抗的能力在他们能力之内，而他们自愿同意，那么我们不得不通过他们的教导找出，为什么在至恶中有如此大的善事，为什么在至善中有这种恶，除非既不是他们所怀疑的是恶，也不是他们用迷信所歪曲的是至善。\par

% structure: p0038 source=h2 target=section reason=relative-heading-rank; base=h2

\section{第13章——从对恶与善的深思可得出不应持两类灵魂之说。即便承认有一类灵魂引诱人作可耻之事，也不能推定这些灵魂本性邪恶，而其他灵魂是至善。}

19. 但倘若我曾教导，或至少自己学到，他们关于那两类灵魂发狂出错的，为什么我以后就该认为他们配得听或就任何事情咨询他们呢？为的是我因此学到，这两类灵魂被指出来，在深思过程中，同意有时站在恶的一边，有时站在善的一边？为什么这倒不是同一个灵魂的标记，这个灵魂凭自由意志能被带到这里那里，被摇动到这里那里呢？因为我自己的经验是感受到我是一个，考虑恶和善，选择其一或另一，但大部分时候一个喜欢，另一个合适，置于两者中间我们摇摆。这也并不奇怪，因为我们现被这样构成，以致通过肉身我们能受感官快乐影响，通过精神受可敬考虑影响。那么我岂非被迫承认两个灵魂？不，我们更好而且远少困难地认出两类善物，其中没有一类是远离\Person{天主}为其作者的，一个灵魂从不同方向、较低和较高，或更正确地说，外在与内在被作用。这就是我们刚才以可感觉的和可理解的名称所考虑的两类，现在我们更愿意更熟悉地称之为属肉体的和属灵的。但对我们来说，禁制肉性之物已经变得困难，因为我们最真实的食粮是属灵的。因为如今我们以大劳力吃这个食粮。因为没有为悖逆之罪的惩罚，我们从未从不朽变为可朽的。所以发生，当我们追求更好的事物时，因与肉身连接所形成的习惯和我们的罪在某种程度上开始攻击我们并拦阻我们的路，有些愚拙的人以最迟钝的迷信怀疑有另一类不属于\Person{天主}的灵魂。\par

20. 然而，即使向他们承认，我们是被另一种低级的灵魂诱骗去行可耻之事，他们也并未由此证明，那些引诱者在本性上是恶的，或被引诱者是最善的。因为有可能，前者出于自己的意志，追求不合法的事，即犯罪，从而从善变为恶；并且他们也可能重新成为善的，但方式是在长时间内停留在罪中，并通过某种隐秘的劝说，将其他灵魂诱骗到自己这边。再者，他们可能并非绝对邪恶，而是在其自身种类中，尽管较低级，却可以在毫无罪恶的情况下行使自身的功能。而那些优越的灵魂，正义（万物的引导者）为其指派了远更卓越的活动，若它们想要追随和模仿那些较低级的灵魂，就成为恶的，不是因为它们模仿邪恶的灵魂，而是因为它们以邪恶的方式模仿。邪恶的灵魂所行的是适合其自身的事，而善良的灵魂所追求的是与其自身相异的事。因此，前者停留在自身的等级中，后者则陷入更低级的等级。正如人模仿野兽时的情况。四足的马行走优美，但若人四肢着地模仿它，谁还会认为他配得上哪怕是草料作为食物？因此，我们通常正确不赞同模仿者，却赞同他所模仿的对象。但我们不赞同，并非因为他没有成功，而是因为他竟然希望成功。因为在马身上，我们赞同那种品质，而正因我们更看重人，故对模仿低等造物感到不悦。同样，在人间，无论传令官如何出色地发出声音，难道元老院议员若比传令官更清晰、更好地发出声音，就为疯狂了吗？取天体为例：月亮照耀时受到赞扬，其运行和变化令关注此事者颇为愉悦。但若太阳想要模仿月亮（我们可以想象它有此类欲望），谁不会大感不悦且理应不悦呢？从这些例子中，我希望明白：即使存在一些灵魂（目前暂且存疑），它们并非因罪恶，而是因本性而致力于身体的功能，并且即使它们与我们有着某种内在的亲和关系，无论它们多么低级，也不应仅仅因为我们在追随它们和热爱物质事物方面是邪恶的，就被视为邪恶。因为我们犯罪是由于热爱物质事物，因为公义要求我们，并且我们本性上能够热爱精神事物，而当我们这样做时，我们在自己种类中就是最好和最幸福的。\par

21. 因此，思虑被剧烈地推向两个方向，时而倾向于犯罪，时而趋向正当行为，这提供了什么证明，迫使我们接受两种灵魂，其中一种的本性来自天主，另一种则不然？当我们能够猜测如此多其他思想交替状态的原因时，这又能证明什么？然而，这些事是晦涩的，而且昏花的头脑徒劳地探究它们，这是任何善于判断事物的人都能看到的。因此，不如那些关于意志和罪所说的事，那些事，我说的是，至高的公义不允许任何运用理性的人对其无知，那些事，若从我们身上除去，就没有任何地方能够开始德行的训练，也没有任何地方能够从恶习的死亡中兴起，那些事，我说，经过足够清晰和透彻的反复思考，使我们确信\OriginalTerm{摩尼派}{Manichaeans}的异端是虚假的。\par

% structure: p0042 source=h2 target=section reason=relative-heading-rank; base=h2

\section{第十四章——再次从悔改的效用证明灵魂并非在本性上为恶。如此确凿的论证唯有因习惯犯错而遭到反驳}

22. 与我刚才关于悔改所说的类似。因为所有心智健全的人都同意——\OriginalTerm{摩尼派}{Manichaeans}本身不仅承认，而且教导——悔改罪过是有益的。我此刻为何要搜集散见于各页的圣经证言呢？这也是本性的呼声；这件事的注意点没有逃过任何愚人的眼睛。如果这不是深深植根于我们的本性，我们就要遭殃了。有人可能说他不犯罪；但没有野蛮人会胆敢说，若犯罪，就不应为此悔改。既然如此，我问，悔改属于两种灵魂中的哪一种？我当然知道它既不属于行恶的人，也不属于不能行善的人。因此，为了使用\OriginalTerm{摩尼派}{Manichaeans}的言辞，若黑暗的灵魂悔改罪过，它就不属于至恶的实体；若光明的灵魂悔改，它就不属于至善的实体；这种有益的悔改态度同时见证了悔改者行过恶，并能行善。那么，若我行事不慎，怎能说从我身上没有恶？或若我未曾如此行，又怎能正确悔改？听另一部分：若我内有善的意志，怎能说从我身上没有善？或若没有善的意志，又怎能正确悔改？因此，他们要么否认悔改有巨大效用，从而被逐出基督教之名并逐出任何为他们观点辩护的想象性论证；要么停止说和教导有两种灵魂，一种毫无恶，另一种毫无善；因为整个派系就依靠这种双头——或更确切地说，是鲁莽的——灵魂多样性来支撑。\par

23. 对我而言，仅仅知道摩尼教徒错了，知道罪必须悔改，这就足够了；然而若我现在按友谊的权利向一位仍认为他们值得听从的朋友进言，并对他说：\Quote{难道你不知道一个人犯了罪后悔改是有益的吗？}他会毫不迟疑地发誓说他明白。那么，若我说服你相信摩尼教是假的，你还会渴望更多吗？让他回答在这件事上他还能渴望什么。很好，到此为止。但当我开始展示那些如同用金钢石锁链捆绑在此命题上（如俗语所说）的确实而必要的论据，并将整场推翻该教派的过程引至结论时，他或许会否认自己知道悔改的益处——而这一点，无论学者还是文盲，无人不知——反而会在我们犹豫不决并仔细思量时，坚持说我们体内有两个灵魂，各自为问题的不同部分提供其特有的帮助。啊，犯罪的习惯！啊，随之而来的罪的惩罚！那时，你使我远离了对如此明显之事的思考，但你伤害了我，因为我未能辨别。但现在，在那些仍不能辨别的至交之中，你因我能辨别而刺痛并折磨我。\par

% structure: p0045 source=h2 target=section reason=relative-heading-rank; base=h2

\section{第十五章——他为曾与自己一同陷入错误的朋友们祈祷}

24. 请留意这些，我恳求你们，亲爱的朋友们。我对你们的心态十分了解。如果你们现在向我承认任何人的心智与理性，那么这些事比我们似乎曾学习或更确切地被迫相信的事要确定得多。伟大的天主，全能的天主，至善的天主，祢是应当被相信和被认识为不可侵犯和不变更的。\OriginalTerm{天主圣三}{Trinal Unity}，天主教会所敬拜的，我，一个在祢仁慈中经验过自身的人，恳求祢，求祢不许那些自幼与我在各种关系中极为和谐相处的人，在敬拜祢的事情上与我分离。我看到在此处尤其要预料到的是，我应当要么在当时就捍卫受到摩尼教徒攻击的天主教圣经，假如我像所说的那样谨慎的话；要么现在应当表明它们可以受到捍卫。但在其他卷帙中，天主将助我如愿，因为我认为这一卷的适度长度已经请求被宽恕。\par

\SourceRecord{Translated by Albert H. Newman.；Nicene and Post-Nicene Fathers, First Series；Vol. 4.；Edited by Philip Schaff.；Buffalo, NY: Christian Literature Publishing Co.,；1887.；<http://www.newadvent.org/fathers/1403.htm>.}
